\subsubsection{Label Ambiguity}\label{sec:label-ambiguity}

The previous challenge was \textbf{dimensionality}: images contain too many input variables (pixels) for a model to learn from. The next challenge is a different kind of problem: \textbf{even if we had a perfect representation of the image, the correct label may not be unique}.

\highspace
\begin{flushleft}
    \textcolor{Green3}{\faIcon{question-circle} \textbf{What is label ambiguity?}}
\end{flushleft}
Suppose we show the following picture: \emph{A man sitting at a restaurant with a beer and sausages on the table}. What is the correct label? Man, Beer, Dinner, Restaurant, Sausages, Food, Table, or something else? All of these describe the same image, and \textbf{none of them is objectively wrong}.

\highspace
So, \definition{Label Ambiguity}, in image classification, means that a \hl{single image can reasonably correspond to more than one possible class, making it impossible to assign a unique ``\emph{correct}'' label based solely on the visual information.}

\highspace
\textcolor{Red2}{\faIcon{question-circle} \textbf{Why is label ambiguity a problem?}} In supervised learning we assume that every training image has \textbf{one ground-truth label}. For example, with the previous image, the dataset label might be ``Restaurant''. However, the network may instead predict ``Dinner''. \emph{Is that prediction incorrect}? From a human point of view, it is not. The image certainly predicts a dinner scene. But from the dataset's point of view, the network is wrong because the annotation only accepts ``Restaurant'' as the correct label. This creates an intrinsic ambiguity in the training data.

\highspace
Also, notice that the ambiguity does not come from the neural network itself (the \textbf{issue is not the network}). It comes from \textbf{the way datasets are labeled}. A single image often contains multiple objects, multiple activities, and multiple concepts. Yet, many datasets force us to choose \textbf{exactly one label}.

\highspace
\textcolor{Red2}{\faIcon{exclamation-triangle} \textbf{Why does this make learning harder?}} During training, the network minimizes the loss assuming there is \textbf{one correct answer}. For example, suppose the dataset says ``Label $=$ Restaurant'', but the network predicts ``Dinner''. Although the prediction is perfectly reasonable, the \hl{loss function still treats it as 100\% wrong.} Consequently, the \textbf{model is penalized for making a semantically valid prediction}.

\highspace
\textcolor{Green3}{\faIcon{check-circle} \textbf{If the problem is a dataset labeling issue, why not just assign multiple labels to each image?}} This is a good idea. In fact, many modern datasets do exactly this. Instead of assigning a single class, they use \textbf{multi-label classification}, where one image can belong to several classes simultaneously. For example, the image of a man at a restaurant could be labeled as both ``Restaurant'' and ``Dinner''. However, throughout this part of the course we are studying the simpler \textbf{single-label classification problem}, where \hl{every image must belong to exactly one class.} Label ambiguity is therefore an \textbf{inherent limitation of this formulation}.